\documentclass[11pt,a4paper]{article}
\usepackage[utf8]{inputenc}
\usepackage[T1]{fontenc}
\usepackage[french]{babel}
\usepackage{tikz}
\usepackage{geometry}
\geometry{margin=2cm}

\begin{document}

\section*{Exercice : Graphe orienté et matrice d'adjacence}

On considère un réseau informatique modélisé par un graphe orienté $G$ comportant 4 sommets notés $A$, $B$, $C$ et $D$.

\subsection*{Représentation graphique du graphe orienté}

\begin{center}
\begin{tikzpicture}[scale=0.9, every node/.style={circle, draw, thick, minimum size=8mm}]
    % Positionnement des sommets
    \node (A) at (0,2) {$A$};
    \node (B) at (3,2) {$B$};
    \node (C) at (3,-1) {$C$};
    \node (D) at (0,-1) {$D$};

    % Arêtes orientées (flèches)
    \draw[->, thick, >=stealth] (A) to[bend left=15] (B);
    \draw[->, thick, >=stealth] (A) to (C);
    \draw[->, thick, >=stealth] (B) to (C);
    \draw[->, thick, >=stealth] (B) to[bend left=15] (D);
    \draw[->, thick, >=stealth] (C) to (A);
    \draw[->, thick, >=stealth] (D) to (B);
    \draw[->, thick, >=stealth] (D) to (A);
\end{tikzpicture}
\end{center}

\end{document}
